\section{Triển khai hệ thống}

\subsection{Kiến trúc triển khai}
Hệ thống được triển khai theo mô hình client-server tách biệt frontend và
backend, đóng gói bằng Docker để đảm bảo tính nhất quán giữa môi trường
phát triển và production:

\begin{itemize}
    \item \textbf{Backend API:} Đóng gói Docker image, triển khai trên
    Render (free/hobby tier phù hợp với ngân sách đồ án sinh viên).
    \item \textbf{Frontend Web Admin:} Triển khai trên Vercel, tận dụng
    CDN toàn cầu giúp cải thiện thời gian tải trang.
    \item \textbf{Cơ sở dữ liệu:} PostgreSQL managed service (Render
    Postgres hoặc Neon), Redis managed instance cho hàng đợi đồng bộ.
    \item \textbf{Mobile App:} Build APK/IPA thử nghiệm qua Expo, phân
    phối nội bộ cho hội đồng chấm và giảng viên hướng dẫn kiểm thử (chưa
    phát hành chính thức lên App Store/Google Play trong phạm vi đồ án).
\end{itemize}

\subsection{Quy trình CI/CD}
Áp dụng GitHub Actions với 2 pipeline:
\begin{itemize}
    \item \textbf{Pipeline Build \& Test:} Kích hoạt khi có Pull Request
    vào nhánh \texttt{develop} — tự động chạy Unit Test và Lint trước
    khi cho phép merge.
    \item \textbf{Pipeline Deploy:} Kích hoạt khi merge vào nhánh
    \texttt{main} — tự động build Docker image, deploy backend lên
    Render và trigger deploy frontend trên Vercel.
\end{itemize}

\subsection{Hướng dẫn cài đặt}
Xem chi tiết đầy đủ tại mục "Installation Guides" — hướng dẫn cài đặt môi
trường phát triển local cho cả Backend (NodeJS/NestJS, Docker) và
Frontend (ReactJS, Node.js 20) được trình bày kèm theo Release Package
bàn giao cuối kỳ.